% SHARD-ID: CA-62-CATEGORICAL-BU-PIPELINE-VISION
% SHARD-TITLE: The Categorical Borchers-Uhlmann Pipeline: a Programme Refinement
% SHARD-SUMMARY: Recasts the indefinite-particle construction as a proposed observable *-algebra-first pipeline rather than a Hilbert-space-first construction.
% SHARD-SUMMARY: Orders category data, free local-word algebra, categorical relations, state selection, GNS quotient, and OAR/Koo-Saleur continuum limits.
% SHARD-SUMMARY: Records the Jones no-go and exact-relation scan failure as load-bearing constraints on the programme.
% SHARD-KEYWORDS: Borchers-Uhlmann, GNS, tensor algebra, anyon chains, Koo-Saleur, OAR, defects, functoriality

\section{The Categorical Borchers--Uhlmann Pipeline: a Programme Refinement}
\label{sec:categorical-bu-pipeline-vision}

\subsection{Status, Sources, and Dependencies}

\textbf{Status: proposal/question.}  This shard records a programme refinement,
not a theorem.  Its central proposal is that the CA-02/CA-03
indefinite-particle grammar should be reinterpreted first as an observable
\(*\)-algebra of local words, and only later as a Hilbert space after a state
and GNS quotient have been chosen.

Dependencies: CA-01, CA-02, CA-03, CA-09, CA-43, CA-44, CA-57, CA-60,
CA-61; CONVENTIONS.md (q); companion shards CA-63, CA-64, CA-65, CA-66
(this block).

% Source: report/sections/03_indefinite_particle_fock.tex:72--96 -- full Fock
%   space as Hilbert direct sum and Hilbert completion of the tensor algebra.
% Source: references/qft/FewsterRejzner2019/source/AQFTIMPRS-Oct2019.tex:724 --
%   field algebra specified by generators Phi(f) over test functions.
% Source: references/qft/FewsterRejzner2019/source/AQFTIMPRS-Oct2019.tex:1901--1913 --
%   free unital *-algebra, two-sided *-ideal of relations, quotient A(M).
% Source: literature/SYNTHESIS.md:8--18 -- categorical variable-N space
%   \oplus Mor(X_c, X_{a_1}\otimes...\otimes X_{a_N}) and pair moves.
% Source: report/sections/09_fibonacci_kinematic_sector.tex:65--93 --
%   fusion admissibility as a diagonal sector projector.
% Source: references/text/HollandsAnyonicChainsAlphaInduction.txt:15--35 --
%   finite anyonic-chain MPO algebra isomorphic to CFT defect algebra; scaling
%   limit identification is conjectural.
% Source: references/text/HollandsAnyonicChainsAlphaInduction.txt:1884--1890 --
%   DHR endomorphisms for Virasoro intervals form a modular tensor category.
% Source: references/text/HollandsAnyonicChainsAlphaInduction.txt:2016--2035 --
%   irreducible defects correspond to central projections in the CFT defect
%   algebra.
% Source: references/text/HollandsAnyonicChainsAlphaInduction.txt:4399--4402 --
%   analytic scaling-limit persistence of the defect algebra is an open question.
% Source: report/sections/43_qubit_vacuum_moment_constraints.tex:77--98 --
%   GNS relation constraints \omega(X^* R Y)=0 on finite probe spaces.
% Source: report/sections/44_qubit_sdp_exclusion_hierarchy.tex:75--94 --
%   infeasibility excludes only the named generator route; feasibility proves
%   neither state existence nor continuum symmetry.
% Source: literature/md/2010.11121/2010.11121.md:126--164 -- OAR pipeline:
%   scaling maps, inductive-limit algebra, projective-limit state, GNS Hilbert
%   representation, and dynamics question.
% Source: literature/md/2010.11121/2010.11121.md:178--180 -- projectively
%   consistent states induce compatible GNS Hilbert-space systems.
% Source: literature/md/2306.16063/2306.16063.md:1094--1124 -- compatible
%   representation nets from projective nets of states via GNS.
% Source: references/text/CFTFromLatticeFermions.txt:221--253 -- Theorem A:
%   Koo-Saleur approximants converge to continuum Virasoro generators for the
%   free-fermion scaling limit.
% Source: references/text/CFTFromLatticeFermions.txt:262--312 -- Theorems B/C:
%   convergence of fermion and Virasoro correlation functions.
% Source: literature/md/2201.11562/2201.11562.md:7 -- braiding gives a real-
%   space RG and the Ising-chain scaling limit gives the Ising CFT vacuum.
% Source: literature/md/2201.11562/2201.11562.md:48--56 -- anyonic-chain input
%   from a fusion category and F-symbol local projectors.
% Source: references/lattice-symmetry/Jones2017NoGo/source/nogo.tex:61--69 --
%   no-go abstract: the block-spin/Temperley-Lieb Hilbert space supports no
%   chiral CFT whose Hamiltonian generates rotation as a limit of lattice
%   rotations.
% Source: references/lattice-symmetry/Jones2017NoGo/source/nogo.tex:951--955 --
%   the main no-go theorem for triadic limits of annular TL representations.
% Source: literature/md/2306.16063/2306.16063.md:1025--1057 -- Thompson/Jones
%   implementers fail for nonzero central charge; Koo-Saleur route succeeds.
% Source: references/lattice-symmetry/ZiniWang2018/source/main.tex:588--590 --
%   scaling limit of quantum theories as a double colimit.
% Source: references/lattice-symmetry/ZiniWang2018/source/main.tex:871--882 --
%   Conjecture 4.3: unitary minimal model VOAs as strong scaling limits with
%   approximating lattice Virasoro observables.

\subsection{Refined North Star}

\textbf{Proposal.}  The 2026-07-05 refinement of the north star is a recipe, as
functorial as possible, which starts from a unitary fusion category or modular
tensor category \(\Cc\), together with any conformal data required by the target
CFT, and produces a family of microscopic lattice models whose scaling limit
realises a chiral or full CFT whose defects and excitations recover the input
category.

\textbf{Proposal.}  The key reinterpretation is algebra-first.  CA-03 builds
\[
  \Fock(H)=\hoplus_{n\ge0}H^{\hotimes n}
\]
as a Hilbert completion and explicitly identifies the underlying tensor
algebra \(T(H)=\bigoplus_n H^{\otimes n}\).  In this shard, the tensor algebra
is not primarily the pre-Hilbert finite-particle subspace.  It is the proposed
free local-word \(*\)-algebra from which observables are generated.  The local
anchor for this style is the free unital \(*\)-algebra of generators
\(\Phi(f)\) over test functions, quotiented by a two-sided \(*\)-ideal of
relations; the registered source builds exactly this object (it does not use
the name Borchers--Uhlmann, which we keep as informal terminology).

\textbf{Proposal.}  Candidate lattice models emerge only after three extra
choices: a reference state \(\omega\), approximate discrete Virasoro or
Koo--Saleur-type residual relations, and quotienting by the state-null
relations.  In symbols, the intended order is
\[
  \text{local-word } *\text{-algebra}
  \quad\stackrel{\omega,\mathcal R}{\longrightarrow}\quad
  \mathfrak N_{\omega,\mathcal R}
  \quad\stackrel{\mathrm{GNS}}{\longrightarrow}\quad
  (\mathcal H_\omega,\pi_\omega,\Omega_\omega),
\]
not ``choose a Hilbert space first, then search for relations''.

\subsection{Five Stages}

\begin{description}
\item[Stage 0: input data.]
\textbf{Status: prior-shard question.}  The input is a unitary fusion category
or modular tensor category, plus conformal/OPE data when the target CFT needs
more than \(\Cc\).  CA-01 states this as the global programme map.  CA-09 gives
the first checked kinematic example: Fibonacci fusion admissibility produces a
sector projector, but full anyons still need coherent \(F/R\)-data.

\item[Stage 1: free algebra of anyonic words.]
\textbf{Status: proposal.}  Replace the Hilbert-first reading of CA-03 by an
observable-algebra reading: generate formal local anyonic words from the
categorical variable-\(N\) sectors
\[
  \bigoplus_{N,c,\mathrm{configs}}
  \operatorname{Mor}(X_c,X_{a_1}\otimes\cdots\otimes X_{a_N}).
\]
The precise algebra generators, involution, locality grading, and completion
policy are fixed in CA-65 and CA-66 (companion shards, this block).

\item[Stage 2: categorical relations.]
\textbf{Status: mixed checked/proposal.}  Fusion admissibility is already a
sector-projector relation in CA-09.  Additional categorical relations should
include defect/MPO relations.  Hollands supplies the finite-size anchor: the
bipartite anyonic-chain MPO algebra is isomorphic to the CFT defect algebra at
finite size, while the persistence of that algebra in the scaling limit is
explicitly left as an analytic question.

\item[Stage 3: reference state and relaxed residuals.]
\textbf{Status: proposal, currently exclusion-mode.}  The state must be chosen
before a Hilbert space is claimed.  CA-43 and CA-44 provide the present finite
architecture: translation-invariant moments, positivity, and residual relation
constraints of the form
\[
  \omega(X^*RY)=0.
\]
In the GNS reading, these are the finite probe shadows of the null-relation
condition.  CA-63 (companion shard, this block) is the construction-mode state
existence route; CA-64 (companion shard, this block) is the relaxation route.

\item[Stage 4: GNS quotient and lattice model.]
\textbf{Status: sourced scaffold plus proposal.}  Once \((\Alg,\omega)\) and a
chosen residual set \(\mathcal R\) are fixed, the candidate Hilbert space and
operator representation are produced by the GNS quotient.  The local OAR
sources support this ordering: build scaling maps and an inductive-limit
algebra, construct compatible/projective limit states, then perform GNS to get
Hilbert-space representations and compatible representation systems.

\item[Stage 5: continuum limit and fidelity.]
\textbf{Status: precedent plus question.}  The continuum-limit technology is
anchored by OAR and the free-fermion Koo--Saleur theorem: lattice Virasoro
approximants and correlation functions converge in that sourced free-fermion
setting.  The anyonic precedent is Stottmeister's braiding-RG note, where
braiding defines a real-space RG for anyonic chains and the Ising-chain example
reaches the Ising CFT vacuum.  The general shape of the open problem is now
also locally sourced: Zini--Wang define scaling limits of quantum theories as
double colimits and conjecture that every unitary minimal model VOA arises as a
strong scaling limit with approximating lattice Virasoro observables.  The
categorical fidelity target remains a question:
\[
  \text{DHR/defect content of the limit} \simeq \Cc?
\]
Hollands makes this a precise open problem by proving a finite algebraic
defect correspondence and asking whether it persists under scaling limits.
\end{description}

\subsection{Two Load-Bearing Constraints}

\textbf{Constraint 1: the Jones no-go changes the order.}  The locally
registered no-go theorem states that the block-spin/Temperley--Lieb
semicontinuous Hilbert space does not support a chiral CFT whose Hamiltonian
generates circle translation as a continuum limit of lattice rotations; in the
main theorem, triadic-rotation-rigid vectors in the triadic limit of an annular
TL representation are necessarily zero.  The OAR discussion corroborates the
programme reading: nontrivial Thompson implementers are not implementable in
the presence of nonzero central charge, and the Koo--Saleur strategy succeeds
in those situations.  The consequence is load-bearing: do not expect a fixed
inductive-limit Hilbert space with dyadic rotations to carry the continuum
geometry first.  The state comes first; geometry and symmetry are recovered
through state-dependent scaling limits and generator convergence.

\textbf{Constraint 2: exact UV relations are too strict.}  CA-57 defines the
current, conservation, and exact boost-witness gates.  CA-60 reports the
checked 99-point scan: no point reaches \texttt{not\_excluded\_algebraic}, and
the locally sourced self-dual TFIM point fails the exact boost witness.  CA-61
and CONVENTIONS.md (q) record the route-scoped reading: this rejects the exact
first-moment route, not the Ising scaling limit.  Therefore Stage 3 must use
quantified state-level or GNS-norm residuals, not exact coefficient identities,
as the default relation language.

\subsection{Functoriality and Fidelity Assessment}

\textbf{Proposal.}  Stages 1 and 2 should be functorial in principle: free
algebras and imposed categorical relations are natural constructions once the
generator and convention data are fixed.  This is not yet a theorem; CA-65
records the intended functor on pairs of a category and a site object.

\textbf{Proposal with sourced scaffold.}  GNS should be treated functorially on
compatible \((\Alg,\omega)\)-pairs.  The local source support is not an
abstract category theorem; it is the OAR architecture in which projective nets
of states induce compatible nets of GNS representations and Hilbert-space
connecting maps.

\textbf{Question.}  The state assignment is the non-functorial crux.  The
realistic target is canonicity by a universal property: define a compact or
closed convex set of states satisfying the relaxed residual constraints, then
prove non-emptiness, uniqueness or controlled faces, and fidelity of the
resulting limit.  The strongest intended theorem would say that this state set
exists and its continuum DHR/defect category recovers the input \(\Cc\).

\subsection{Boundary}

This shard does not fix the categorical local-word algebra; the \(F/R\)-symbol
gauge, object indexing, bracketing, and locality/completion conventions are
fixed in CA-65/CA-66, not here.

It does not prove existence of the reference state.  CA-43/CA-44 are
exclusion-mode finite constraints; CA-63 is the companion construction-mode
shard.

It does not fix the relaxation convention.  The exact-relation scan failure
only says exact UV boost witnesses are too strict for the current route; CA-64
must define the norm, scale, tolerance, and limiting quantifiers.

It does not prove a new continuum theorem.  The only continuum precedents cited
here are the sourced free-scalar/OAR, free-fermion Koo--Saleur, and Ising
braiding-RG examples; the Zini--Wang minimal-model statement is cited as a
conjecture, not a result.

It does not claim that SDP feasibility or finite residual smallness implies
symmetry.  CA-44 already records that feasibility does not prove a global
state, positive energy, or a continuum Poincare/Virasoro representation.
